% \section*{Résumé}
\begin{resume}
    Les avions de combat modernes ont de plus en plus besoin d'un niveau de furtivité au moins modéré et la conception de la forme doit supporter une partie de cette contrainte. Cependant, la fréquence élevée des radars à courte portée rend les méthodes de calcul de la section transversale radar à haute fidélité, telles que la méthode des éléments de frontière, trop coûteuses pour être utilisées dans un processus d'optimisation sans gradient. D'autre part, les méthodes asymptotiques ne permettent pas de prédire avec précision la RCS de formes complexes telles que les cavités d'admission. Il est donc nécessaire de trouver des méthodes efficaces et précises pour calculer le gradient des méthodes de calcul de la section transversale radar haute fidélité en fonction des paramètres de la forme.
    Dans cet article, nous proposons une formulation adjointe pour la méthode des éléments de frontière afin de calculer efficacement ces gradients. Nous présentons une nouvelle approche pour calculer le gradient des calculs de section transversale radar haute fidélité à l'aide de la méthode des éléments de frontière. Notre méthode utilise une formulation adjointe qui permet un calcul efficace de ces gradients. Cette méthode est particulièrement utile dans les problèmes d'optimisation de la forme où des méthodes précises et efficaces sont cruciales pour concevoir des avions de combat modernes dotés de capacités de furtivité. En évitant la nécessité de résoudre le problème adjoint dans certains cas, notre formulation fournit une solution plus rationnelle tout en conservant une grande précision. Nous démontrons l'efficacité de notre méthode en effectuant une optimisation de forme sur différentes formes, y compris des géométries simples comme la sphère
\end{resume}